\documentclass[10pt]{article}
\usepackage[letterpaper,margin=0.72in]{geometry}
\usepackage{amsmath,amssymb}
\usepackage{booktabs,longtable,array}
\usepackage[colorlinks=true,urlcolor=blue,linkcolor=black]{hyperref}
\setlength{\parindent}{0pt}
\setlength{\parskip}{0.45em}
\renewcommand{\arraystretch}{1.12}
\title{\vspace{-1.5em}SAEE lb120: Post-submission Operational Clarification}
\author{Zhang Bin}
\date{Local support artifact; not submitted to the portal}

\begin{document}
\maketitle

\begin{abstract}
This document clarifies the operational meaning and provenance of quantities
reported in the ALIFE 2026 Late-Breaking Abstract submission \texttt{lb120}.
It introduces no new experiment, result, theory, runtime, kernel, candidate
law, or GSP semantic change. The submitted PDF remains frozen. The evidence is
limited to one 100-generation long-horizon artifact set and one six-generation
Phase II analysis artifact set; no cross-run universality, formal attractor
existence, open-ended evolution, benchmark superiority, external validation,
publication, proceedings inclusion, DOI assignment, or completed conference
registration is claimed. Linklings records the submission as accepted and the
author confirmation as completed.
\end{abstract}

\section{Submission and venue boundary}

The portal state observed on 2 July 2026 was \texttt{Under Evaluation} for
submission \texttt{lb120}. On 18 July 2026, the Linklings decision and author
confirmation state was observed as \texttt{Accept (Confirmed)}. This records
acceptance and confirmed intent to present; it does not establish publication,
proceedings inclusion, DOI assignment, completed conference registration, or
external validation. The official ALIFE 2026 call states that Late-Breaking Abstracts are at most two
pages excluding references, may report new ideas and work in progress, are
reviewed for relevance and quality, are not included in proceedings, and are
presented as posters if accepted. Submissions are non-anonymous and the review
process is single-blind.\footnote{\url{https://2026.alife.org/call-for-papers/}}

The frozen submitted PDF has one page. Its SHA-256 is:
\begin{quote}\small\ttfamily
aef09e556b2e91b2374b51371164d887\\
e2db40c39a260b3dac2ef7772d15ece8
\end{quote}
The submitted source and PDF are not modified by this support artifact.

\section{Evidence surfaces}

Two surfaces are compressed in the LBA and must remain distinct.

\begin{enumerate}
\item The v1.0 long-horizon surface records 100 generations with configured
population size 8, collapse checks, fitness-variance tendency, genome-ID
turnover, and a lineage DAG.
\item The Phase II surface analyzes six generations from a v0.8 record and
emits semantic-drift values, categorical state signatures, empirical attractor
candidates, and pointwise regime labels.
\end{enumerate}

The later label \texttt{stable\_lineage\_basin} is a local phase-diagram
compression of both surfaces, not the literal attractor identifier emitted by
Phase II and not a formal estimate of a basin in continuous phase space.

\section{Long-horizon operational definitions}

\subsection{Configuration and determinism}

The recorded configuration is:
\begin{quote}\small
\path{generation_count=100}; \path{population_size=8};
\path{deterministic_seed=enabled}; \path{logging_level=full_trace}; and
\path{seed_path=kernel/examples/seed_genome.json}.
\end{quote}
The deterministic-seed field is a mode flag, not a numeric random seed. v1.0
child identifiers are derived from truncated SHA-256 digests and mutations are
deterministic weight shifts.

\subsection{Collapse and fitness-variance tendency}

A collapse event is recorded for generation $t$ when
\[
 |P_t| < 2
 \quad\text{or}\quad
 \mathrm{survivor\_count}_t = 0.
\]
No such event appears in the recorded run. This is a local absence under the
implemented rule, not a proof that collapse is impossible.

Let $F_{\mathrm{first}}$ be the mean of per-generation fitness variance over
the first ten generations and $F_{\mathrm{last}}$ the corresponding mean over
the final ten. The analyzer emits \texttt{converging} when
$F_{\mathrm{last}} < 0.75F_{\mathrm{first}}$, \texttt{diverging} when
$F_{\mathrm{last}} > 1.25F_{\mathrm{first}}$, and
\texttt{stable\_variance} otherwise. The current artifact is labeled
\texttt{converging}; this is not mathematical limit convergence.

\subsection{Turnover and lineage accounting}

For consecutive genome-ID sets $P_{t-1}$ and $P_t$, population turnover is the
Jaccard distance
\[
 d_J(P_{t-1},P_t)
 = 1-\frac{|P_{t-1}\cap P_t|}{\max(1,|P_{t-1}\cup P_t|)}.
\]
The first generation is assigned zero. The reported mean is \texttt{0.545531}.

The values \texttt{808/1590} mean 808 nodes and 1590 directed edges in one
declared \texttt{lineage\_dag}. They do not mean retained nodes versus raw
events. Branching density is
\[
 \frac{1590}{808}=1.967822.
\]
The flag \texttt{lineage\_integrity\_preserved=true} requires the declared
graph type \texttt{lineage\_dag} and verifies that every edge endpoint occurs
in the node-ID set. The reporting function does not independently perform a
topological cycle test.

\section{Phase II operational definitions}

\subsection{Semantic drift}

Let $A$ be the set of lowercased identity-anchor reference terms and $D_t$ the
set of lowercased dominant terms after any bounding intervention. The recorded
drift is
\[
 \delta_t=1-\frac{|D_t\cap A|}{\max(1,|A|)}.
\]
It is a term-overlap distance, not an embedding-space semantic distance. The
bound is \texttt{0.32}; the six-generation artifact reports
$\delta_t=0.2$ at every point and therefore \texttt{semantic\_drift\_max=0.2}.

\newpage
\begin{samepage}
\subsection{State signature and empirical attractor candidate}

Each trajectory point receives four categorical components:

{\small
\begin{center}
\begin{tabular}{@{}>{\raggedright\arraybackslash}p{0.18\linewidth}p{0.74\linewidth}@{}}
\toprule
Component & Implemented rule \\
\midrule
Identity & \texttt{stable} iff no identity break and mean identity score is at least 0.58; otherwise \texttt{fragile}. \\
Drift & \texttt{bounded} iff semantic drift after control is at most 0.32; otherwise \texttt{unbounded}. \\
Population & \texttt{dense} iff population count is at least 10; otherwise \texttt{sparse}. \\
Mutation & \texttt{active} iff explanation-driven mutation count is positive; otherwise \texttt{quiet}. \\
\bottomrule
\end{tabular}
\end{center}
}

An exact signature occurring at least twice is emitted as an empirical
attractor candidate. The current signature is
\texttt{identity:stable|drift:bounded|population:sparse|mutation:active}, with
support at generations 1--6. There is no sliding window, continuous distance,
clustering algorithm, dominance ratio, dispersion threshold, or recurrence-gap
threshold. Hence \emph{attractor} means a recurring discrete signature in this
record, not a formal continuous-state attractor proof.
\end{samepage}

\subsection{Regime rule}

The point classifier applies rules in order: (1) \texttt{chaotic\_regime} if
an identity break occurs or drift exceeds 0.55; (2) \texttt{collapse\_regime}
if population is at most 2; (3) \texttt{stable\_regime} if the point supports
an attractor candidate and drift is at most 0.32; (4)
\texttt{exploratory\_regime} if mutation-event count is at least 4; and (5)
\texttt{stable\_regime} otherwise.

Rule order matters: the six current points have at least five mutation events,
but satisfy the stable-attractor rule first. All six are labeled
\texttt{stable\_regime}. The classifier reports zero label changes. A later
projection counts five adjacent self-transitions, giving local ratio 5/5; this
is not a general transition-probability estimate.

\section{Related-work boundary}

Tierra and Avida establish the core digital-evolution context, with Avida in
particular emphasizing controlled protocols and extensive analysis tools
\cite{ray1992,ofria2004}. MODES proposes comparable metrics for hallmarks of
open-ended dynamics \cite{dolson2019}; lb120 implements none of those metrics
and makes no open-ended-evolution claim. Ancestry-based analysis work develops
lineage and phylogeny metrics \cite{dolson2020}; lb120's node and edge counts
are not phylogenetic richness, divergence, or equivalent measures.
Reproducibility guidance for evolutionary computation motivates recording
code, data, factors, and environment while distinguishing fixed-instance
claims from broader inference \cite{lopez2021}.

The narrow contribution is therefore a repository-linked observation protocol
for one frozen object: it exposes how identity-bounded feedback, a lineage DAG,
categorical regimes, and local evidence compression map to inspectable
artifacts. It is not a platform-breadth or algorithm-performance claim.

\newpage
\section{Provenance and known gaps}

\begin{center}
\begin{tabular}{@{}p{0.31\linewidth}p{0.61\linewidth}@{}}
\toprule
Field & Recorded value or status \\
\midrule
Long-horizon config & \texttt{saee\_experiments/configs/experiment\_config.yaml} \\
Long-horizon trace & \texttt{saee\_experiments/output/demo-run/evolution\_trace.jsonl} \\
Lineage report & \texttt{saee\_experiments/reports/lineage\_statistics.json} \\
Phase II attractor map & \texttt{saee\_phase2/output/demo-run/attractor\_map.json} \\
Phase II regime log & \texttt{saee\_phase2/output/demo-run/regime\_transition\_log.json} \\
Original execution commit & Not recorded in current run artifacts \\
Exact original command & Not recorded \\
Numeric random seed & Not applicable as a recorded field; deterministic mode flag only \\
Original OS/Python and lockfile & Not recorded \\
Original CPU/GPU & Not recorded \\
\bottomrule
\end{tabular}
\end{center}

The repository HEAD observed during this documentation pass was
\texttt{f6ac41f4b068377e7778e8c3d83b99bd8382debc}. It is a current
documentation snapshot and is not represented as the original execution
commit.

\section{Claim-safe conclusion}

The strongest reconstructable statement is that the recorded 100-generation
v1.0 surface remained viable under its collapse rule and preserved declared-DAG
endpoint integrity, while the six-generation Phase II surface repeated one
bounded-drift categorical signature and received six stable-regime labels.
This is local observational evidence only.

\begin{thebibliography}{9}
\bibitem{ray1992}
T. S. Ray. An Approach to the Synthesis of Life. In \emph{Artificial Life II}, 1992.
\url{https://tomray.me/pubs/alife2/Ray1991AnApproachToTheSynthesisOfLife.pdf}

\bibitem{ofria2004}
C. Ofria and C. O. Wilke. Avida: A Software Platform for Research in
Computational Evolutionary Biology. \emph{Artificial Life}, 10(2):191--229,
2004. doi:10.1162/106454604773563612.

\bibitem{dolson2019}
E. L. Dolson, A. E. Vostinar, M. J. Wiser, and C. Ofria. The MODES Toolbox:
Measurements of Open-Ended Dynamics in Evolving Systems. \emph{Artificial
Life}, 25(1):50--73, 2019. doi:10.1162/artl\_a\_00280.

\bibitem{dolson2020}
E. Dolson, A. Lalejini, S. Jorgensen, and C. Ofria. Interpreting the Tape of
Life: Ancestry-Based Analyses Provide Insights and Intuition about Evolutionary
Dynamics. \emph{Artificial Life}, 26(1):58--79, 2020.
doi:10.1162/artl\_a\_00313.

\bibitem{lopez2021}
M. L\'opez-Ib\'a\~nez, J. Branke, and L. Paquete. Reproducibility in
Evolutionary Computation. \emph{ACM Transactions on Evolutionary Learning and
Optimization}, 1(4):1--21, 2021. doi:10.1145/3466624.
\end{thebibliography}

\end{document}
